% =====================================================================
\section{Methods}
\label{sec:methods}
% =====================================================================

\subsection{Bayesian inference}
\label{sec:bayes}

Every strategy runs the same loop, Fig.~\ref{fig:bayes-loop}: it chooses a drive frequency $d_n$, takes one
measurement $y_n$ there, updates the posterior over the line parameters with it, and stops when the stop rule of
Sec.~\ref{sec:locators} fires. The strategies differ only in the first step; the posterior and its update, described
here, are shared. The unknowns are those of Sec.~\ref{sec:fisher}, the parameter of interest $\theta$ and the
nuisances $\phi$, together with the noise level $\sigma$.

\begin{figure}[t]
\centering
% Flow chart of the measurement loop shared by every strategy (Sec. Bayesian inference, Fig. fig:bayes-loop).
% Single column. Only the "choose" box (shaded) differs between strategies.
\begin{tikzpicture}[
  font=\small,
  node distance=4.5mm,
  box/.style={draw, rounded corners=2pt, align=center, text width=4.4cm, inner sep=3pt},
  pick/.style={box, fill=gray!20, thick},
  test/.style={draw, diamond, aspect=2.6, align=center, inner sep=1pt, font=\footnotesize},
  arr/.style={-{Stealth[length=2mm]}},
]
\node[box] (prior) {Prior: particles $\{(\theta_i,\phi_i)\}$ drawn\\ over the search window};
\node[pick, below=of prior] (choose) {Choose drive frequency $d_n$\\[1pt]
  {\footnotesize sweep: next grid point\\ SBED: largest EIG in the focus}};
\node[box, below=of choose] (measure) {Measure $y_n$ at $d_n$};
\node[box, below=of measure] (update) {Reweight particles $w_i$\\[1pt]
  {\footnotesize Student-$t$ likelihood, $\sigma$ integrated out}};
\node[test, below=of update] (ess) {weight on few\\ particles?};
\node[box, right=3mm of ess, text width=1.6cm, font=\footnotesize] (resample) {resample +\\ Liu--West jitter};
\node[test, below=of ess] (stop) {stop rule\\ met?};
\node[box, below=of stop] (report) {Report $\hat\theta_N$ and $\hat\sigma_N$, Eq.~\eqref{eq:sigma-hat}};

\draw[arr] (prior) -- (choose);
\draw[arr] (choose) -- (measure);
\draw[arr] (measure) -- (update);
\draw[arr] (update) -- (ess);
\draw[arr] (ess) -- node[above, font=\scriptsize] {yes} (resample);
\draw[arr] (ess) -- node[right, font=\scriptsize] {no} (stop);
\draw[arr] (resample) |- (stop);
\draw[arr] (stop) -- node[right, font=\scriptsize] {yes} (report);
\coordinate (back) at ([xshift=-3mm]choose.west);
\draw[arr] (stop.west) -- node[above, font=\scriptsize, pos=0.3] {no} (stop.west -| back) -- (back) -- (choose.west);
\end{tikzpicture}

\caption{The measurement loop shared by every strategy. Only the choice of the next drive frequency (shaded)
differs between strategies; the particle posterior, its update and the stop rule are the same for all of them.
The noise level $\sigma$ is integrated out in each particle, so the update weighs by a Student-$t$ likelihood.}
\label{fig:bayes-loop}
\end{figure}

The posterior over $(\theta,\phi)$ is a weighted cloud of particles $\{(\theta_i,\phi_i),w_i\}_{i=1}^{N_p}$, updated
by sequential Monte Carlo \cite{doucet2001smc}. Each particle is one hypothesis for the line. It holds the Zeeman
shift $\delta_{Z,i}$, so $\theta_i=D-\delta_{Z,i}$ with $D$ fixed (the simulations measure the mirror
transition $D+\delta_{Z,i}$, which is equivalent by symmetry), and the nuisances: the width of the line and its
contrast $C$. The mechanisms that set the width, Eq.~\eqref{eq:linewidth}, are not inferred; only the width they
produce is. The initial cloud is drawn from the prior.

The noise is Gaussian, Eq.~\eqref{eq:noise}, but its level $\sigma$ is not known in advance, since technical noise
adds to the shot noise by an amount that depends on the setup. Each particle therefore carries a closed-form
posterior on $\sigma^2$, given its line, and $\sigma^2$ is integrated out \cite{storvik2002,carvalho2010pl}: a
measurement reweights each particle by a Student-$t$ likelihood instead of a Gaussian with a plug-in $\sigma$, which
would drive the estimate of $\sigma$ to its lower bound. When the weight has concentrated on a few particles, the
cloud is resampled \cite{kitagawa1996} and jittered with the Liu--West kernel \cite{liuwest2001}, which preserves its
mean and covariance. The noise model, the resampling safeguards and all settings are in \SM{sm:algorithm}.

The particle filter reports the posterior mean and the weighted particle standard deviation,
\begin{equation}
  \hat\theta_N=\sum_i w_i \theta_i,\qquad
  \hat\sigma_N=\Big[\sum_i w_i(\theta_i-\hat\theta_N)^2\Big]^{1/2},
  \label{eq:sigma-hat}
\end{equation}
the claimed uncertainty whose meaning Sec.~\ref{sec:protocol} discusses.

\subsection{Measurement strategies}
\label{sec:locators}

We compare Bayesian SBED with the strategy an experimentalist would use without adaptive design, a sweep, run at
several grid spacings; the best spacing is the yardstick SBED is measured against. All strategies infer with the
particle filter of Sec.~\ref{sec:bayes}, so the comparison isolates the choice of measurements from the inference,
and $\hat\sigma_N$ means the same for all of them. All run under the same stop rule and the same measurement budget:
a run stops once the claimed uncertainty of every parameter has stayed below its ceiling for several consecutive
measurements and the dip test below finds a dip.

\paragraph{Dip test.} The test asks where the measurements, not the particles, show a dip. Each measurement is scored
against a baseline taken from the upper part of the data, in units of the noise from the noise posterior. Contiguous
stretches of measured drive frequencies, about as wide as a plausible line, get a boxcar score, and a stretch is
significant when its score passes a family-wise false-alarm threshold over all stretches tested. Each connected run of
significant stretches is one dip.

\paragraph{Sweep.} SimpleSweep measures a uniform grid over the window, a few points per FWHM of the
prior-mean line, and repeats whole passes; the spacing is the one setting a sweep has, and several are run.

\paragraph{Bayesian SBED.} SBED is the particle filter choosing each $d$ by Eq.~\eqref{eq:eig-var}, maximized over
candidate drive frequencies concentrated on the steep sides of the posterior-mean dip and on the dips found. The
candidates are restricted to the focus: one section per detected dip, padded by one half width on each side, which
widens at once when a new dip appears and narrows only after a delay. A strategy that has piled its measurements onto
a suspected dip cannot be contradicted elsewhere, so with a small probability, largest while much of the window is
unmeasured, SBED instead measures the middle of the widest unmeasured gap, a falsification probe.

\subsection{Simulation setup and evaluation}
\label{sec:protocol}

\paragraph{What the claimed uncertainty is.} The particle filter reports $\hat\sigma_N$, Eq.~\eqref{eq:sigma-hat}: a moment of the
posterior, defined at every $N$. Its run average is compared with the bound from the first measurement on
(Sec.~\ref{sec:budget}). What fails before the posterior is Gaussian
(Sec.~\ref{sec:eig}) is the usual reading of $\hat\sigma_N$: with a uniform window or several live
candidates it is not a 68\% interval, and the posterior entropy is below $\tfrac12\ln(2\pi e\,\hat\sigma_N^2)$, so
information gains, per-step gains as variance ratios, and coverage at Gaussian quantiles are only meaningful in the
data regime. We therefore (i) compare $\hat\sigma_N$ with the bound as a run average at all $N$, (ii) test coverage,
the true error against $\hat\sigma_N$, only once the posterior is Gaussian, or with quantiles of the particle
distribution itself, and (iii) keep this separate from miscalibration of the particle filter, which makes
$\hat\sigma_N$ too small at any $N$.
% Code: smc_marginal.py uncertainty() = weighted particle std; entropy() assumes a Gaussian (logdet of the
% covariance), so it overstates the entropy before the posterior is Gaussian.

\paragraph{Use cases.} A use case is one simulated signal at one noise level, set by the SNR, the drive window, the
Gaussian share of the width and the isotope (\SM{sm:algorithm}). Every metric is computed within one use
case, over repeats whose true signal is the same for every strategy, so strategies are compared run by run; several
use cases are summarized by the median over them.

\paragraph{Metrics.} The accuracy at the end of a run is not a quality metric: every run stops once its claimed
uncertainty is below the ceiling, so its final error mostly reflects the ceiling. We report instead whether the
claimed uncertainty is honest, through its coverage, the share of runs whose true error lies within $\hat\sigma_N$
and within $2\hat\sigma_N$; how close each strategy comes to the bound of Eq.~\eqref{eq:vantrees}, through the ratio
of the bound to the mean squared true error at each $N$; and the cost, the number of measurements after which the
true error stays below a target. Each run is also read against the bound through its information rate, the growth of
$\hat\sigma_N^{-2}$ per measurement in units of $I_{\rm max}$, Eq.~\eqref{eq:istar}; a rate above 1 means the claimed
uncertainty fell faster than the data allow. Measurements to stop are compared only between strategies whose claimed
uncertainty is equally honest.

\subsection{Experimental setup}
\label{sec:setup}

We measure the lower transition $f_-$ only: in our setup it needs less microwave power than $f_+$, which keeps the
heating of the stage low and the amplifier within its linear range.
\tbd{Sample, optics, microwave delivery, detection and the per-shot integration time $t_s$.}
